\section{Methods}
\label{sec:methodology}

\subsection{Outcome Definitions}
\label{sec:meth_outcomes}

For each trip $i$ with onboard speed vector $\mathbf{v}_i = (v_{i1}, \ldots, v_{iT_i})$ recorded at approximately 10-second intervals, we compute eight primary behavioral outcomes organized into three categories:

\textit{Safety.} Harsh acceleration count: the number of consecutive-point acceleration events exceeding 0.5~m/s$^2$, where acceleration is computed as $a_{it} = (v_{i,t+1} - v_{it}) \times (1000/3600) / 10$. The threshold is calibrated for the $\sim$10-second GPS interval (a standard 2~m/s$^2$ threshold yields near-zero events at this sampling rate). Harsh deceleration count is computed analogously with $a_{it} < -0.5$~m/s$^2$. Speed coefficient of variation: $\text{CV}_i = \sigma(\mathbf{v}_i) / \bar{v}_i$.

\textit{Demand.} Trip distance (meters) and travel time (seconds) are recorded directly by the platform. At the city-month level, trip count and active user count serve as demand-side indicators.

\textit{Riding dynamics.} Cruise fraction: the proportion of speed readings within $\pm$3~km/h of the trip mean speed, capturing steady-state riding. Zero-speed fraction: the proportion of speed readings equal to zero, capturing stop-and-go behavior.

Mean speed, maximum speed, and 85th percentile speed are computed as \textit{manipulation checks}: they confirm the governor's mechanical effect but are not treated as primary findings, as speed reduction under a speed cap is tautological.

\subsection{Cross-Sectional Behavioral Profiles}
\label{sec:meth_gee}

We model each primary outcome as a function of operating mode and controls via generalized estimating equations (GEE) to account for within-user correlation \citep{zeger1986longitudinal}:
\begin{equation}
    g\big(E[Y_i]\big) = \beta_0 + \beta_1 \text{STD}_i + \beta_2 \text{ECO}_i + \boldsymbol{x}_i^{\top} \boldsymbol{\gamma}
    \label{eq:gee}
\end{equation}
where TUB is the reference mode, $g(\cdot)$ is the link function (log link with Poisson family for count outcomes; identity link with Gaussian family for continuous outcomes), and $\boldsymbol{x}_i$ includes age group, time of day, day type (weekday/weekend), city fixed effects, fraction of trip on major roads, and log-transformed trip distance. We specify exchangeable working correlation and cluster by user ID. Robust (sandwich) standard errors are used throughout. Separate GEE models are estimated for each of the eight outcomes plus the manipulation check (mean speed). A random subsample of 200,000 trips (reservoir sampling) is used for computational feasibility.

Within-subject paired comparisons are conducted for users who rode in multiple modes. For each user who took trips in both TUB and STD (or TUB and ECO), we compute the user-level mean of each outcome in each mode and report the paired difference with Cohen's $d$ effect sizes.

\subsubsection{Individual fixed effects}

To address potential endogeneity from rider self-selection into modes, we supplement the GEE with individual fixed effects (FE) models that absorb all time-invariant user characteristics (age, risk preference, riding skill):
\begin{equation}
    Y_{it} = \alpha_i + \beta_1 \text{STD}_{it} + \beta_2 \text{ECO}_{it} + \boldsymbol{x}_{it}^{\top} \boldsymbol{\gamma} + \epsilon_{it}
    \label{eq:fe}
\end{equation}
where $\alpha_i$ is a user fixed effect, and $\boldsymbol{x}_{it}$ includes time of day, day type, road class (major road, cycling infrastructure), and log-transformed trip distance. Continuous outcomes are estimated via fixed-effects OLS (\texttt{feols}) and count outcomes via fixed-effects Poisson (\texttt{fepois}), implemented in \texttt{pyfixest} \citep{fischer2024pyfixest}. Standard errors are clustered at the user level (CRV1). A 2-million-trip subsample is used for computational feasibility; after removing singleton users (those with only one trip), approximately 1.8 million trips from 343{,}000 users remain.

Comparing GEE and FE estimates provides a Hausman-style endogeneity diagnostic. If GEE and FE coefficients agree, between-user variation is exogenous and GEE is efficient. If they diverge---particularly if GEE yields larger mode effects---user self-selection inflates the cross-sectional estimates, and the FE within-user estimates are preferred for causal interpretation.

\subsection{Feature Importance: LightGBM with SHAP}
\label{sec:meth_shap}

To assess relative feature importance without distributional assumptions, we train separate LightGBM models for three outcome categories: safety (binary: any harsh event), dynamics (cruise fraction, continuous), and safety (speed CV, continuous). Each model uses a 300,000-trip subsample with features including mode indicators, log-distance, road class, age group, time of day, and day type. SHAP (SHapley Additive exPlanations) values \citep{lundberg2017unified} are computed on a 5,000-trip subset to quantify each feature's contribution, providing model-agnostic rankings of feature importance across outcome categories.

\subsection{Multi-Outcome Difference-in-Differences}
\label{sec:meth_did}

The December 2023 TUB ban provides causal identification. We construct a city-month panel (52 cities $\times$ 11 months, February--December 2023) and estimate a two-way fixed effects (TWFE) DiD for each outcome separately:
\begin{equation}
    Y_{ct}^{(j)} = \alpha_c + \gamma_t + \beta^{(j)} \cdot (\text{Post}_t \times \text{TUB}^{\text{Nov}}_c) + \epsilon_{ct}
    \label{eq:twfe}
\end{equation}
where $Y_{ct}^{(j)}$ is the mean of outcome $j$ in city $c$ at month $t$, $\alpha_c$ and $\gamma_t$ are city and month fixed effects, $\text{Post}_t = \mathbf{1}(t = \text{Dec~2023})$, and $\text{TUB}^{\text{Nov}}_c$ is the November 2023 TUB mode share in city $c$ (continuous treatment intensity). Standard errors are clustered by city. The parameter $\beta^{(j)}$ estimates the causal effect of governor imposition on outcome $j$.

We estimate Equation~\ref{eq:twfe} for all six primary outcomes plus two demand outcomes (trip count, active user count), yielding eight tests. To control the family-wise error rate, we apply Bonferroni correction with $\alpha_{\text{Bonf}} = 0.05/8 = 0.00625$.

Parallel trends are assessed via an event study specification:
\begin{equation}
    Y_{ct}^{(j)} = \alpha_c + \gamma_t + \sum_{k \neq \text{Nov}} \delta_k^{(j)} \cdot \mathbf{1}(t = k) \cdot \text{TUB}^{\text{Nov}}_c + \epsilon_{ct}
    \label{eq:event_study}
\end{equation}
with November as the reference period. Dose-response is estimated by regressing city-level outcome changes (November $\rightarrow$ December) on TUB share changes, with bootstrapped 95\% confidence intervals (1,000 replications). A placebo test applies the TWFE specification using September 2023 as a fictitious treatment date, restricting to February--November.

\subsection{Compensation Test: Within-User Mode-Switcher Design}
\label{sec:meth_compensation}

To distinguish behavioral compensation from pool composition, we conduct a within-user mode-switcher analysis. ``Switchers'' are users who used TUB in October--November 2023 (pre-ban) and rode STD/ECO in December (post-ban). ``Never-TUB'' users are those who never used TUB in any month. For each user in both groups, we compute the mean of each outcome in STD/ECO trips during the pre-period (October--November) and post-period (December), then estimate a user-level DiD:
\begin{equation}
    \text{DiD}^{(j)} = \overline{\Delta Y}_{\text{switcher}}^{(j)} - \overline{\Delta Y}_{\text{never-TUB}}^{(j)}
    \label{eq:switcher_did}
\end{equation}
where $\Delta Y$ is the within-user pre-post change. Significance is assessed via Welch's $t$-test on the deltas; effect sizes are reported as Cohen's $d$. A user-level placebo test using August--September $\rightarrow$ October (no treatment) validates the design.

The key distinction: if the STD/ECO-only DiD (Section~\ref{sec:meth_did}) shows significant effects but the within-user switcher DiD shows negligible Cohen's $d$, the aggregate change is attributable to \textit{composition} (former TUB users joining the STD/ECO pool are inherently different riders) rather than \textit{compensation} (individual riders adjusting behavior).

\subsection{Behavioral Escalation: Experience Mediation and Survival}
\label{sec:meth_escalation}

We analyze how riding behavior evolves with experience. For each user, trips are ranked chronologically and grouped into experience bins ($1, 2, 3, 4{-}5, 6{-}10, \ldots, 101{-}200$ trips). At each bin, we compute: (a)~TUB adoption rate, (b)~mean of each outcome across all modes, and (c)~mean of each outcome in STD/ECO modes only. The gap between (b) and (c) quantifies the TUB-mediated component of the experience effect. We report mediation percentages for each outcome.

Multi-endpoint survival analysis models time (in trips) to three first events: (1)~first harsh event (any harsh acceleration or deceleration), (2)~first high-CV trip (speed CV $>$ 75th percentile), and (3)~first speeding trip (max speed $>$ 25~km/h). Kaplan-Meier curves stratified by TUB usage (ever-TUB vs.\ never-TUB) are compared via the log-rank test. Cox proportional hazards models estimate hazard ratios for TUB usage indicators (ever-TUB indicator and TUB usage fraction).
